\section{Equipo Scrum y Eventos}
\subsection{Accountabilities del equipo}
La Scrum Team est\'a compuesta por Product Owner, Scrum Master y Developers. Es un equipo peque\~no, multifuncional y auto-gestionado \cite{schwaber2020}.

\begin{table}[htbp]
\centering
\renewcommand{\arraystretch}{1.15}
\begin{tabularx}{\textwidth}{>{\bfseries}p{3.1cm} X}
\toprule
Rol & Responsabilidad principal \\
\midrule
Product Owner & Maximiza el valor del producto y ordena el Product Backlog. \\
Scrum Master & Facilita Scrum, remueve impedimentos y protege el timebox. \\
Developers & Construyen el incremento y planifican el trabajo del Sprint. \\
\bottomrule
\end{tabularx}
\end{table}

\subsection{Criterios de trabajo}
\begin{itemize}[leftmargin=*]
    \item Un Sprint, un objetivo.
    \item Las decisiones t\'ecnicas pertenecen a Developers.
    \item El Product Owner prioriza por valor.
\end{itemize}

\subsection{Eventos Scrum}
\begin{table}[htbp]
\centering
\renewcommand{\arraystretch}{1.08}
\begin{tabularx}{\textwidth}{>{\bfseries}p{3cm} p{3cm} X}
\toprule
Evento & Timebox oficial & Uso en el proyecto \\
\midrule
Sprint & Un mes o menos & Sprint de 2 semanas. \\
Sprint Planning & Hasta 8 horas / mes & Planificar alcance, objetivo y trabajo. \\
Daily Scrum & 15 minutos & Sincronizar avance y bloqueos. \\
Sprint Review & Hasta 4 horas / mes & Revisar incremento y feedback. \\
Sprint Retrospective & Hasta 3 horas / mes & Definir mejoras del proceso. \\
\bottomrule
\end{tabularx}
\end{table}

\subsection{Prop\'osito operativo}
Los eventos sirven para inspeccionar y adaptar el trabajo, no para generar reuniones innecesarias \cite{schwaber2020}.
